\documentclass[11pt,a4paper]{article}

\usepackage[T1]{fontenc}
\usepackage[utf8]{inputenc}
\usepackage[brazil]{babel}
\usepackage{lmodern}
\usepackage{microtype}
\usepackage[margin=2cm]{geometry}
\usepackage{hyperref}
\usepackage{titlesec}
\usepackage{xcolor}

\hypersetup{
  colorlinks=true,
  linkcolor=black,
  urlcolor=black,
  citecolor=black,
  pdfborder={0 0 0}
}

\definecolor{sectioncolor}{HTML}{111111}
\definecolor{accentcolor}{HTML}{1A1A1A}

\pagestyle{empty}
\setlength{\parindent}{0pt}
\setlength{\parskip}{5pt}
\linespread{1.05}
\selectfont

\titleformat{\section}
  {\bfseries\large\color{sectioncolor}}
  {}{0pt}{}[\vspace{-2pt}{\color{sectioncolor}\titlerule[0.8pt]}]
\titlespacing*{\section}{0pt}{12pt}{8pt}

\newcommand{\cvrole}[4]{%
  \vspace{4pt}%
  {\bfseries\color{accentcolor}#1}\hfill{\itshape #2}\\[-2pt]%
  {\itshape #3}\hfill #4\\[6pt]%
}

\begin{document}

\begin{center}
  {\fontsize{20}{24}\selectfont\bfseries Kamille Hartmann Maccagnan}\\[8pt]
  {\large Governança de TI \textbar{} PMO \textbar{} Power BI \& Indicadores}\\[10pt]
  Campo Comprido, Curitiba-PR \enspace\textbullet\enspace (41) 9 9928-4424 \enspace\textbullet\enspace
  \href{mailto:kamillehartmannm@gmail.com}{kamillehartmannm@gmail.com}\\[3pt]
  \href{https://linkedin.com/in/kamille-hartmann-maccagnan/}{linkedin.com/in/kamille-hartmann-maccagnan/}
\end{center}

\vspace{6pt}

\section{Resumo Profissional}

Profissional com formação em Gestão da Tecnologia da Informação e Medicina Veterinária, com experiência em governança de dados, PMO, padronização de processos e monitoramento de indicadores em ambientes corporativos e industriais. Atua na documentação de processos, elaboração de relatórios executivos e dashboards em Power BI, e interface entre áreas de TI, operações e stakeholders. Vivência em acompanhamento de portfólio de projetos, análise de riscos, gestão de demandas e melhoria contínua, com perfil analítico, organizado e comunicativo. Conhecimento em metodologias de gestão de projetos e práticas de governança, com familiaridade em ambientes que exigem integração entre tecnologia e operação.

\section{Experiência Profissional}

\cvrole{Anjuss}{Analista de Suporte de TI}{07/2025 -- Atual}

Atuo na padronização de processos e governança operacional de sistemas, com criação e atualização de documentações, POPs para a rede de lojas e POPs voltados à equipe de T.I., assegurando conformidade e consistência na execução das rotinas. Lidero treinamentos e capacitação de novos funcionários, promovendo adoção padronizada de procedimentos e disseminação de conhecimento entre as áreas. Desenvolvo dashboards em Power BI para monitoramento de indicadores operacionais e de desempenho, produzindo relatórios que apoiam a gestão e a tomada de decisão. Organizo demandas internas, identifico oportunidades de melhoria de processos e implemento automações com Power Automate para aumentar eficiência operacional e confiabilidade das informações.

\cvrole{HORSE do Brasil}{Estagiária de TI (IS/IT LATAM)}{07/2024 -- 07/2025}

Atuei em frentes de PMO, governança de dados e suporte em ambiente multinacional industrial, com acompanhamento de portfólio de projetos, cronogramas, prazos e entregas. Elaborei relatórios executivos e dashboards para monitoramento de performance de projetos e iniciativas estratégicas, apoiando lideranças com indicadores para tomada de decisão. Atuei na interface entre áreas de TI, operações de chão de fábrica, fornecedores e stakeholders internacionais, promovendo alinhamento e comunicação entre tecnologia e operação. Utilizei o ServiceNow para gestão e acompanhamento de demandas e chamados, e apliquei o Power Automate em iniciativas de automação de processos. Participei de documentação de processos, análise de riscos, melhoria contínua e evolução de práticas de governança de projetos e dados.

\cvrole{Restaurante Familiar}{Analista de Dados e Suporte Técnico}{01/2023 -- 07/2024}

Atuei na análise de performance operacional e resultados financeiros para suporte à gestão do negócio. Estruturei do zero a base de dados em Excel e dashboards em Power BI para acompanhamento de indicadores de desempenho, vendas e resultados, com padronização de informações para maior confiabilidade dos dados. Realizei análises de desempenho, identifiquei desvios e oportunidades de melhoria em processos operacionais, apoiando decisões de otimização de recursos. Atuei no alinhamento entre áreas operacionais e administrativas, com foco em eficiência, controle de indicadores e suporte analítico à liderança.

\section{Competências Técnicas}

\textbf{Governança e Processos:} Governança de dados, padronização de processos, documentação operacional, POPs, mapeamento de processos, melhoria contínua, gestão de demandas, análise de riscos, Boas Práticas PMOBOK, Scrum, Kanban.

\textbf{Indicadores e Relatórios:} Monitoramento de KPIs, relatórios executivos e gerenciais, dashboards em Power BI (DAX, Power Query), análise de dados, Excel, SQL, acompanhamento de performance operacional e financeira.

\textbf{Gestão de Projetos e Stakeholders:} PMO, acompanhamento de portfólio de projetos, controle de cronogramas e entregas, interface entre TI, operação, fornecedores e lideranças, comunicação multilíngue, gestão de relacionamento.

\textbf{Ferramentas:} Power BI, ServiceNow, Power Automate, Google Sheets, Trello, Monday, Spotfire, Pacote Office.

\textbf{Idiomas:} Inglês Avançado \enspace\textbullet\enspace Espanhol Básico \enspace\textbullet\enspace Coreano Básico

\section{Projetos e Conquistas}

1º lugar no Transformation Day Renault 2023 com solução orientada a dados e melhoria de processos. Participação no Hackathon Bosch 2023 com foco em resolução de problemas de negócio. Desenvolvimento de dashboards executivos para monitoramento de indicadores de desempenho. Estruturação de KPIs para análise de performance operacional. Automação e padronização de fluxos de dados e relatórios com Power Automate integrado ao Power BI.

\section{Formação Acadêmica}

\textbf{PUCPR} \hfill Curitiba, PR\\
Graduação em Gestão da Tecnologia da Informação \hfill Previsão de conclusão: 06/2026\\[6pt]

\textbf{Escola Conquer} \hfill Curitiba, PR\\
Pós-graduação em Liderança e Gestão em Tecnologia \hfill Conclusão: 2024\\[6pt]

\textbf{Universidade Tuiuti do Paraná} \hfill Curitiba, PR\\
Graduação em Medicina Veterinária \hfill 06/2019\\[6pt]

\textbf{Qualittas} \hfill Curitiba, PR\\
Pós-graduação em Clínica Médica e Cirúrgica de Animais Exóticos e Pets Não Convencionais \hfill 12/2022

\end{document}
